%%%%%%%%%%%%%%%%%%%%%%%%%%%%%%%%%%%%%%%%%%%%%%%%%%%%%%%%%%%%%
\section{模型评价与讨论}

\subsection{方法学特色}

本文的方法论以三项原则区别于单纯的"跑出一个数"：

\textbf{判据先行、实验拍板。} 全部 10 项建模决策（界面相位、折射率、偏振、衬底模型、背景方案、噪声估计、条纹区界、周期竞争阈值、融合权重方案、一致性检验水平）均在实验前以预注册的定量判据冻结，而后以数据实验定案，判据不可事后修改。这一机制有效防止了"试到好看为止"的数据美化。

\textbf{负结果入预算。} 落选方案并非丢弃，而是将其量化结果归档，作为系统误差区间和模型升级的依据：（1）Drude 自由拟合失败（$\gamma$ 误差棒 $\pm3700\%$，$d$ 漂移 $+31\%$）$\to$ v2.0 固定文献参数路线；（2）参数化色散与 $d$ 简并（$n$ 撞物理下界 2.30）$\to$ 常数 $n$ 定稿 + 系统区间；（3）频域高通背景方案（Gibbs 振铃，条纹幅度畸变 709\%）$\to$ 触发红线出局。负结果的量化归档使系统误差区间有据可查，并直接指向模型升级方向。

\textbf{可辨识性意识。} $n$-$d$ 简并（EXP-2）和 Drude-$d$ 纠缠（EXP-4）两次实证说明"参数越多越准"在此类反问题中不成立。当待估参数与噪声/模型误差耦合时，增加自由参数反而恶化可辨识性——这一教训催生了 v2.0 的"固定文献参数、不增加自由度"路线。

\textbf{选择与验证分离。} 模型选型仅用 10\degree 数据，15\degree 数据独立验证——避免"用同一份数据既选模型又验证模型"的循环论证。

\subsection{模型优点}

\begin{enumerate}[itemindent=2em]
\item \textbf{数学连续性：} v1.0（双光束）$\to$ v1.5（三域融合）$\to$ v2.0（Airy 多光束）是同一物理过程在不同精度层级上的展开，v2.0 在 $q\to0$ 时精确退化到 v1.0（偏差 $3.9\times 10^{-16}$），三版本不是三个独立模型。

\item \textbf{验证体系完备：} 四条独立验证线（闭环、互证、残差、交叉）覆盖了算法自身、数据自洽、模型充分性、路径独立性四个维度——深度超过可查基准。

\item \textbf{不确定度诚实：} 系统误差以区间 $[8.0,8.9]\,\mu$m 报告，明确标注不对称性（$-0.0/+0.9\,\mu$m）；统计 $\sigma$ 下 L2 的 FAIL 如实记录而非掩盖；$\sigma_p$ 双口径张力如实入档。

\item \textbf{框架可推广：} 三域互证 + 两级一致性检验的框架天然输出"模型是否需要升级"的判据（残差结构 = 触发信号），可直接推广至多层结构分析。
\end{enumerate}

\subsection{模型缺点与局限}

\begin{enumerate}[itemindent=2em]
\item \textbf{折射率依赖文献值：} $n_1=2.55$ 取自外部文献（Goldberg 2001），非本实验独立测定。在文献值缺失或不适用的场景（如新材料的未知外延层），折射率包络自反演路线更为鲁棒——但 EXP-2 已证实单角度数据不可辨识 $n$ 与 $d$，需要变角测量系列或独立色散数据。

\item \textbf{衬底模型区间较宽：} $\sigma_p$ 的幅度反演值与跨角度一致值存在张力（700 vs 4600\,cm$^{-1}$），对应 SiC 厚度全区间 $[8.03,9.15]\,\mu$m。收窄该区间需要更精确的衬底色散模型（如 Lorentz 振子显式建模剩余射线带 793.9 和 970.1\,cm$^{-1}$）。

\item \textbf{Si 反演框架降级：} E1/E2 对 Si 形态失效（条纹周期非常数），三域互证在 Si 上退化为单算法 $+$ 双角度，验证强度有所降低。

\item \textbf{过渡层/粗糙度未建模：} v2.0 拟合残差 8--10\%（Si）仍含结构——工艺中的掺杂过渡层和界面微观粗糙度尚属模型的未覆盖物理。

\item \textbf{面内均匀性未验证：} 无重复测量数据，无法评估晶圆不同位置的厚度偏差。
\end{enumerate}